\section{Introducción}
Hasta no hace muchas décadas, lo obtención de energía eléctrica para abastecer el sistema se obtenía mayormente mediante la quema de combustible, un ejemplo de esto es la obtención de electricidad mediante generadoras termoeléctricas, gracias al cambio de paradigma en el que se esta implementando cada vez más la instalación de ERNC (Energías renovables no convencionales)  en el sistema eléctrico nacional, es fundamental entender la importancia y el funcionamiento de estas nuevas tecnologías.\\

Chile es uno de los mejores países del mundo para invertir en paneles fotovoltaicos y turbinas eólicas, gracias a las características climáticas del norte y el sur. [0] Por lo que es clave aprovechar este recurso al máximo e innovar en nuevas tecnologías para este objetivo. \\

El desafío de reducir las emisiones de gases de efecto invernadero es un gran desafío, donde todos los países del mundo deben contribuir, en el año 2017 Chile firmó  el Acuerdo de París [9], su objetivo es limitar el aumento de temperatura por debajo de los $2°C$ respecto a la era preindustrial, e idealmente limitar este aumento a un máximo de $1.5°C$.\\

Entendiendo la importancia de las ERNC como serían los paneles fotovoltaicos y las turbinas eólicas, en este laboratorio se estudiará el comportamiento básico de estas 2 tecnologías, donde se construirán las curvas características de cada tecnología, entendiendo además los principales desafíos en incorporar estas tecnologías a la red eléctrica, la eficiencia y sus fenómenos principales asociados que afectan la magnitud de potencia activa inyectada a la red como serían el efecto Sombra para el panel fotovoltaico, el ángulo y velocidad del viento para la turbina eólica.